\subsection{ARM: \RU{3 аргумента}\EN{3 arguments}\ES{3 argumentos}}

\RU{В ARM традиционно принята такая схема передачи аргументов в функцию: 
4 первых аргумента через регистры \Reg{0}-\Reg{3}; а остальные~--- через стек}
\EN{ARM's traditional scheme for passing arguments (calling convention) behaves as follows:
the first 4 arguments are passed through the \Reg{0}-\Reg{3} registers; the remaining arguments via the stack}.
\ES{El esquema tradicional de ARM para pasar argumentos (convención de llamadas) funciona de la siguiente manera: los primeros 4 argumentos se pasan a través de los registros \Reg{0}-\Reg{3}; los argumentos restantes a través de la pila.}
\RU{Это немного похоже на то, как аргументы передаются в}\EN{This resembles the arguments passing scheme in}\ES{Esto se asemeja al esquema de paso de argumentos en} 
fastcall~(\myref{fastcall}) \OrENRU win64~(\myref{sec:callingconventions_win64}).

\subsubsection{32-\RU{битный}\EN{bit}\ES{bits} ARM}

\myparagraph{\NonOptimizingKeilVI (\ARMMode)}

\begin{lstlisting}[caption=\NonOptimizingKeilVI (\ARMMode)]
.text:00000000 main
.text:00000000 10 40 2D E9   STMFD   SP!, {R4,LR}
.text:00000004 03 30 A0 E3   MOV     R3, #3
.text:00000008 02 20 A0 E3   MOV     R2, #2
.text:0000000C 01 10 A0 E3   MOV     R1, #1
.text:00000010 08 00 8F E2   ADR     R0, aADBDCD     ; "a=%d; b=%d; c=%d"
.text:00000014 06 00 00 EB   BL      __2printf
.text:00000018 00 00 A0 E3   MOV     R0, #0          ; return 0
.text:0000001C 10 80 BD E8   LDMFD   SP!, {R4,PC}
\end{lstlisting}

\RU{Итак, первые 4 аргумента передаются через регистры \Reg{0}-\Reg{3}, по порядку: 
указатель на формат-строку для \printf
в \Reg{0}, затем 1 в \Reg{1}, 2 в \Reg{2} и 3 в \Reg{3}.}
\EN{So, the first 4 arguments are passed via the \Reg{0}-\Reg{3} registers in this order:
a pointer to the \printf format string in 
\Reg{0}, then 1 in \Reg{1}, 2 in \Reg{2} and 3 in \Reg{3}.}
\ES{Entonces, los primeros 4 argumentos se pasan a través de los registros \Reg{0}-\Reg{3} en este orden: un puntero a la cadena de formato de \printf en \Reg{0}, luego 1 en \Reg{1}, 2 en \Reg{2} y 3 en \Reg{3}.}

\RU{Инструкция на}\EN{The instruction at}\ES{La instrucción en} \TT{0x18} \RU{записывает}\EN{writes}\ES{escribe} 0 \RU{в}\EN{to}\ES{en} \Reg{0}%
\EMDASH{}\RU{это выражение в Си}\EN{this is} \IT{return 0}\EN{ C-statement}\ES{esto es la sentencia \IT{return 0} de C}.

\RU{Пока что здесь нет ничего необычного}\EN{There is nothing unusual so far}\ES{No hay nada inusual hasta ahora}.

\OptimizingKeilVI \RU{генерирует точно такой же код}\EN{generates the same code}\ES{genera exactamente el mismo código}.

\myparagraph{\OptimizingKeilVI (\ThumbMode)}

\begin{lstlisting}[caption=\OptimizingKeilVI (\ThumbMode)]
.text:00000000 main
.text:00000000 10 B5        PUSH    {R4,LR}
.text:00000002 03 23        MOVS    R3, #3
.text:00000004 02 22        MOVS    R2, #2
.text:00000006 01 21        MOVS    R1, #1
.text:00000008 02 A0        ADR     R0, aADBDCD     ; "a=%d; b=%d; c=%d"
.text:0000000A 00 F0 0D F8  BL      __2printf
.text:0000000E 00 20        MOVS    R0, #0
.text:00000010 10 BD        POP     {R4,PC}
\end{lstlisting}

\RU{Здесь нет особых отличий от неоптимизированного варианта для режима ARM.}
\EN{There is no significant difference from the non-optimized code for ARM mode.}
\ES{No hay una diferencia significativa con el código no optimizado para el modo ARM.}

\myparagraph{\OptimizingKeilVI (\ARMMode) + \RU{убираем}\EN{let's remove}\ES{eliminemos} return}
\label{ARM_B_to_printf}

\RU{Немного переделаем пример, убрав}\EN{Let's rework example slightly by removing}\ES{Modifiquemos ligeramente el ejemplo eliminando} \IT{return 0}:

\begin{lstlisting}
#include <stdio.h>

void main()
{
	printf("a=%d; b=%d; c=%d", 1, 2, 3);
};
\end{lstlisting}

\RU{Результат получится необычным:}
\EN{The result is somewhat unusual:}
\ES{El resultado es algo inusual:}

\begin{lstlisting}[caption=\OptimizingKeilVI (\ARMMode)]
.text:00000014 main
.text:00000014 03 30 A0 E3   MOV     R3, #3
.text:00000018 02 20 A0 E3   MOV     R2, #2
.text:0000001C 01 10 A0 E3   MOV     R1, #1
.text:00000020 1E 0E 8F E2   ADR     R0, aADBDCD     ; "a=%d; b=%d; c=%d\n"
.text:00000024 CB 18 00 EA   B       __2printf
\end{lstlisting}

\index{ARM!\Registers!Link Register}
\index{ARM!\Instructions!B}
\index{Function epilogue}
\RU{Это оптимизированная версия (\Othree) для режима ARM, и здесь мы видим последнюю инструкцию 
\TT{B} вместо привычной нам \TT{BL}}\EN{This is the optimized (\Othree) version for ARM mode and this time we see \TT{B} as the last instruction instead of the familiar \TT{BL}}.
\ES{Esta es la versión optimizada (\Othree) para el modo ARM, y esta vez vemos \TT{B} como la última instrucción en lugar de la familiar \TT{BL}.}
\RU{Отличия между этой оптимизированной версией и предыдущей, скомпилированной без оптимизации, 
ещё и в том, 
что здесь нет пролога и эпилога функции (инструкций, сохраняющих состояние регистров \TT{\Reg{0}} и \ac{LR})}%
\EN{Another difference between this optimized version and the previous one (compiled without optimization)
is the lack of function prologue and epilogue (instructions preserving the \TT{\Reg{0}} and \ac{LR} registers values)}.
\ES{Otra diferencia entre esta versión optimizada y la anterior (compilada sin optimización) es la falta de prólogo y epílogo de la función (instrucciones que preservan los valores de los registros \TT{\Reg{0}} y \ac{LR}).}
\index{x86!\Instructions!JMP}
\RU{Инструкция \TT{B} просто переходит на другой адрес, без манипуляций с регистром \ac{LR}, то есть
это аналог \JMP в x86}%
\EN{The \TT{B} instruction just jumps to another address, without any manipulation of the \ac{LR} register,
similar to \JMP in x86}.
\ES{La instrucción \TT{B} simplemente salta a otra dirección, sin ninguna manipulación del registro \ac{LR}, similar a \JMP en x86.}
\RU{Почему это работает нормально? Потому что этот код эквивалентен предыдущему.}
\EN{Why does it work? Because this code is, in fact, effectively equivalent to the previous.}
\ES{¿Por qué funciona esto? Porque este código es, de hecho, efectivamente equivalente al anterior.}
\RU{Основных причин две: 1) стек не модифицируется, как и \glslink{stack pointer}{указатель стека} \ac{SP}; 2) вызов функции \printf последний, 
после него ничего не происходит}\EN{There are two main reasons: 1) neither the stack nor \ac{SP} (the \gls{stack pointer}) is modified;
2) the call to \printf is the last instruction, so there is nothing going on afterwards}.
\ES{Hay dos razones principales: 1) ni la pila ni \ac{SP} (el \glslink{stack pointer}{puntero de pila}) se modifican; 2) la llamada a \printf es la última instrucción, por lo que no ocurre nada más después.}
\RU{Функция \printf, отработав, просто возвращает управление по адресу, записанному в \ac{LR}.}
\EN{On completion, the \printf function simply returns the control to the address 
stored in \ac{LR}.}
\ES{Al finalizar, la función \printf simplemente devuelve el control a la dirección almacenada en \ac{LR}.}
\RU{Но в \ac{LR} находится адрес места, откуда была вызвана наша функция!
А следовательно, управление из \printf вернется сразу туда.}
\EN{Since the \ac{LR} currently stores the address of the point from where our function
was called then the control from \printf will be returned to that point.}
\ES{Dado que \ac{LR} actualmente almacena la dirección del punto desde donde se llamó a nuestra función, el control de \printf se devolverá a ese punto.}
\RU{Значит нет нужды сохранять \ac{LR}, потому что нет нужны модифицировать \ac{LR}}%
\EN{Therefore we do not need to save \ac{LR} because we do not need to modify \ac{LR}}.
\ES{Por lo tanto, no necesitamos guardar \ac{LR} porque no necesitamos modificar \ac{LR}.}
\RU{А нет нужды модифицировать \ac{LR}, потому что нет иных вызовов функций, кроме \printf, к тому же, после этого вызова не нужно ничего здесь больше делать}%
\EN{And we do not need to modify \ac{LR} because there are no other function calls except \printf. Furthermore,
after this call we do not to do anything else}!
\ES{Y no necesitamos modificar \ac{LR} porque no hay otras llamadas a funciones excepto \printf. Además, ¡después de esta llamada no necesitamos hacer nada más!}
\RU{Поэтому такая оптимизация возможна}\EN{That is the reason such optimization is possible}.
\ES{Esa es la razón por la que tal optimización es posible.}

\RU{Эта оптимизация часто используется в функциях, где последнее выражение~--- это вызов другой функции.}
\EN{This optimization is often used in functions where the last statement is a call to another function.}
\ES{Esta optimización se utiliza a menudo en funciones donde la última sentencia es una llamada a otra función.}

\RU{Ещё один похожий пример описан здесь}\EN{A similar example is presented here}:
\myref{jump_to_last_printf}.
\ES{Se presenta un ejemplo similar aquí: \myref{jump_to_last_printf}.}

\subsubsection{ARM64}

\myparagraph{\NonOptimizing GCC (Linaro) 4.9}

\lstinputlisting[caption=\NonOptimizing GCC (Linaro) 4.9]{patterns/03_printf/ARM/ARM3_O0.lst.\LANG}

\index{ARM!\Instructions!STP}
\RU{Итак, первая инструкция STP (Store Pair) сохраняет \ac{FP} (X29) and \ac{LR} (X30) в стеке.}
\EN{The first instruction STP (Store Pair) saves \ac{FP} (X29) and \ac{LR} (X30) in the stack.}
\ES{La primera instrucción STP (Store Pair) guarda \ac{FP} (X29) y \ac{LR} (X30) en la pila.}
\RU{Вторая инструкция \TT{ADD X29, SP, 0} формирует стековый фрейм.}
\EN{The second \TT{ADD X29, SP, 0} instruction forms the stack frame.}
\ES{La segunda instrucción \TT{ADD X29, SP, 0} forma el marco de pila (stack frame).}
\RU{Это просто запись значения \ac{SP} в X29.}
\EN{It is just writing the value of \ac{SP} into X29.}
\ES{Es simplemente escribir el valor de \ac{SP} en X29.}

\index{ARM!\Instructions!ADRP/ADD pair}
\RU{Далее уже знакомая пара инструкций \TT{ADRP}/\ADD формирует указатель на строку.}%
\EN{Next, we see the familiar \TT{ADRP}/\ADD instruction pair, which forms a pointer to the string.}
\ES{A continuación, vemos el familiar par de instrucciones \TT{ADRP}/\ADD, que forma un puntero a la cadena.}
\index{ARM64!lo12}
\RU{\IT{lo12} означает младшие 12 бит, т.е., линкер запишет младшие 12 бит адреса метки LC1 в опкод инструкции \ADD.}%
\EN{\IT{lo12} meaning low 12 bits, i.e., linker will write low 12 bits of LC1 address into the opcode of \ADD instruction.}
\ES{\IT{lo12} significa los 12 bits inferiores, es decir, el enlazador escribirá los 12 bits inferiores de la dirección LC1 en el código de operación de la instrucción \ADD.}

\RU{\TT{\%d} в формате \printf это 32-битный \Tint, так что 1, 2 и 3 заносятся в 32-битные части регистров.}
\EN{\TT{\%d} in \printf string format is a 32-bit \Tint, so the 1, 2 and 3 are loaded into 32-bit register parts.}
\ES{\TT{\%d} en el formato de cadena de \printf es un \Tint de 32 bits, por lo que 1, 2 y 3 se cargan en las partes de registro de 32 bits.}

\Optimizing GCC (Linaro) 4.9 \RU{генерирует почти такой же код}\EN{generates the same code}\ES{genera casi el mismo código}.
